\section{AIME I 1999}
        \begin{enumerate}
        	\item (\textit{AIME I 1999}) Find the smallest prime that is the fifth term of an increasing arithmetic sequence, all four preceding terms also being prime.


 

	\item (\textit{AIME I 1999}) Consider the parallelogram with vertices $(10,45),$ $(10,114),$ $(28,153),$ and $(28,84).$  A line through the origin cuts this figure into two congruent polygons.  The slope of the line is $m/n,$ where $m_{}$ and $n_{}$ are relatively prime positive integers.  Find $m+n.$


 

	\item (\textit{AIME I 1999}) Find the sum of all positive integers $n$ for which $n^2-19n+99$ is a perfect square.


 

	\item (\textit{AIME I 1999}) The two squares shown share the same center $O_{}$ and have sides of length 1. The length of $\overline{AB}$ is $\frac{43}{99}$ and the area of octagon $ABCDEFGH$ is $\frac{m}{n}$ where $m_{}$ and $n_{}$ are relatively prime positive integers.  Find $m+n.$


 
                    \begin{figure}[H]
                        \centering
                        \includegraphics[width=0.5\textwidth]{Problem/AIME_1999_I_4_Problem_Figure_1.png}
                    \end{figure}
                    


 

	\item (\textit{AIME I 1999}) For any positive integer $x_{}$, let $S(x)$ be the sum of the digits of $x_{}$, and let $T(x)$ be $|S(x+2)-S(x)|.$  For example, $T(199)=|S(201)-S(199)|=|3-19|=16.$  How many values of $T(x)$ do not exceed 1999?


 

	\item (\textit{AIME I 1999}) A transformation of the first quadrant of the coordinate plane maps each point $(x,y)$ to the point $(\sqrt{x},\sqrt{y}).$  The vertices of quadrilateral $ABCD$ are $A=(900,300), B=(1800,600), C=(600,1800),$ and $D=(300,900).$  Let $k_{}$ be the area of the region enclosed by the image of quadrilateral $ABCD.$  Find the greatest integer that does not exceed $k_{}.$


 

	\item (\textit{AIME I 1999}) There is a set of 1000 switches, each of which has four positions, called $A, B, C$, and $D$.  When the position of any switch changes, it is only from $A$ to $B$, from $B$ to $C$, from $C$ to $D$, or from $D$ to $A$.  Initially each switch is in position $A$.  The switches are labeled with the 1000 different integers $(2^{x})(3^{y})(5^{z})$, where $x, y$, and $z$ take on the values $0, 1, \ldots, 9$.  At step $i$ of a 1000-step process, the $i$-th switch is advanced one step, and so are all the other switches whose labels divide the label on the $i$-th switch.  After step 1000 has been completed, how many switches will be in position $A$?


 

	\item (\textit{AIME I 1999}) Let $\mathcal{T}$ be the set of ordered triples $(x,y,z)$ of nonnegative real numbers that lie in the plane $x+y+z=1.$  Let us say that $(x,y,z)$ supports $(a,b,c)$ when exactly two of the following are true: $x\ge a, y\ge b, z\ge c.$  Let $\mathcal{S}$ consist of those triples in $\mathcal{T}$ that support $\left(\frac 12,\frac 13,\frac 16\right).$  The area of $\mathcal{S}$ divided by the area of $\mathcal{T}$ is $m/n,$  where $m_{}$ and $n_{}$ are relatively prime positive integers, find $m+n.$


 

	\item (\textit{AIME I 1999}) A function $f$ is defined on the complex numbers by $f(z)=(a+bi)z,$ where $a_{}$ and $b_{}$ are positive numbers.  This function has the property that the image of each point in the complex plane is equidistant from that point and the origin.  Given that $|a+bi|=8$ and that $b^2=m/n,$ where $m_{}$ and $n_{}$ are relatively prime positive integers.  Find $m+n.$


 

	\item (\textit{AIME I 1999}) Ten points in the plane are given, with no three collinear.  Four distinct segments joining pairs of these points are chosen at random, all such segments being equally likely.  The probability that some three of the segments form a triangle whose vertices are among the ten given points is $m/n,$ where $m_{}$ and $n_{}$ are relatively prime positive integers.  Find $m+n.$


 

	\item (\textit{AIME I 1999}) Given that $\sum_{k=1}^{35}\sin 5k=\tan \frac mn,$ where angles are measured in degrees, and $m_{}$ and $n_{}$ are relatively prime positive integers that satisfy $\frac mn<90,$ find $m+n.$


 

	\item (\textit{AIME I 1999}) The inscribed circle of triangle $ABC$ is tangent to $\overline{AB}$ at $P_{},$  and its radius is 21.  Given that $AP=23$ and $PB=27,$ find the perimeter of the triangle.


 

	\item (\textit{AIME I 1999}) Forty teams play a tournament in which every team plays every other ($39$ different opponents) team exactly once.  No ties occur, and each team has a $50 \%$ chance of winning any game it plays.  The probability that no two teams win the same number of games is $m/n,$ where $m_{}$ and $n_{}$ are relatively prime positive integers.  Find $\log_2 n.$


 

	\item (\textit{AIME I 1999}) Point $P_{}$ is located inside triangle $ABC$ so that angles $PAB, PBC,$ and $PCA$ are all congruent.  The sides of the triangle have lengths $AB=13, BC=14,$ and $CA=15,$ and the tangent of angle $PAB$ is $m/n,$ where $m_{}$ and $n_{}$ are relatively prime positive integers.  Find $m+n.$


 

	\item (\textit{AIME I 1999}) Consider the paper triangle whose vertices are $(0,0), (34,0),$ and $(16,24).$  The vertices of its midpoint triangle are the midpoints of its sides.  A triangular pyramid is formed by folding the triangle along the sides of its midpoint triangle.  What is the volume of this pyramid?


 

\end{enumerate}